\documentclass[12pt,a4paper]{article}

\usepackage[utf8]{inputenc}
\usepackage[T1]{fontenc}
\usepackage[french]{babel}
\usepackage{geometry}
\usepackage{xcolor}
\usepackage{hyperref}
\usepackage{listings}
\usepackage{longtable}
\usepackage{booktabs}
\usepackage{enumitem}

\geometry{margin=2.3cm}
\setlength{\parskip}{0.6em}
\setlength{\parindent}{0pt}

\hypersetup{
    colorlinks=true,
    linkcolor=blue!60!black,
    urlcolor=blue!60!black
}

\definecolor{codebg}{HTML}{F7F8FA}
\definecolor{codeborder}{HTML}{D8DDE6}
\definecolor{codekw}{HTML}{0F4C81}
\definecolor{codestring}{HTML}{7A3E00}
\definecolor{codecomment}{HTML}{4D7A4D}

\lstdefinestyle{code}{
    backgroundcolor=\color{codebg},
    basicstyle=\ttfamily\small,
    keywordstyle=\color{codekw}\bfseries,
    stringstyle=\color{codestring},
    commentstyle=\color{codecomment},
    frame=single,
    rulecolor=\color{codeborder},
    breaklines=true,
    columns=fullflexible,
    showstringspaces=false,
    tabsize=4
}

\lstset{style=code}

\title{\textbf{Tariqi Legal AI}\\
\large Assistant juridique RAG du code de la route marocain}
\author{Projet Python, RAG et Streamlit}
\date{2026}

\begin{document}

\maketitle
\tableofcontents
\newpage

\section{Résumé}

Tariqi Legal AI est un assistant juridique informatif consacré au code de la route marocain.
Il aide l'utilisateur à comprendre les infractions, les amendes, le permis à points et les
procédures pratiques comme le paiement, la déclaration ou la réclamation.

Le projet ne vise pas à préparer l'examen du permis. Il vise plutôt à fournir une réponse
claire, sourcée et prudente à partir de sources institutionnelles marocaines.

\section{Objectifs}

\begin{itemize}[leftmargin=1.2cm]
    \item Répondre aux questions juridiques routières à partir de sources fiables.
    \item Citer les sources utilisées dans chaque réponse.
    \item Calculer les montants indicatifs et points retirés à partir d'un CSV structuré.
    \item Proposer des procédures simples pour paiement, déclaration et réclamation.
    \item Fournir une application Streamlit utilisable et une interface CLI.
\end{itemize}

\section{Sources}

Le système utilise prioritairement les sources suivantes :

\begin{longtable}{p{4cm} p{8cm} p{2cm}}
\toprule
\textbf{Source} & \textbf{Usage} & \textbf{Niveau} \\
\midrule
\endhead
SGG & Loi n° 52-05, textes consolidés, décrets & A+ \\
Bulletin Officiel & Référence juridique opposable & A+ \\
Ministère du Transport & Décret relatif aux règles de circulation & A+ \\
Khadamat NARSA & Infractions, paiement, déclaration, réclamation & A \\
NARSA & Permis à points et tableaux de retrait de points & A \\
\bottomrule
\end{longtable}

Les blogs, forums, vidéos, posts de réseaux sociaux et PDF non traçables ne sont pas utilisés
comme base principale de réponse.

\section{Architecture}

\begin{lstlisting}
tariqi-legal-ai/
|-- app/                  Interface Streamlit
|-- data/
|   |-- raw/              Sources brutes locales non versionnées
|   |-- seed/             Corpus de départ
|   |-- structured/       CSV infractions et procédures
|-- docs/                 Documentation et rapport
|-- scripts/              Commandes CLI
|-- src/tariqi/           Code métier
|-- tests/                Tests unitaires
|-- vectorstore/          Index généré
\end{lstlisting}

\section{Pipeline RAG}

Le pipeline RAG fonctionne en deux phases.

\subsection{Indexation}

\begin{enumerate}[leftmargin=1.2cm]
    \item Charger les documents.
    \item Nettoyer le texte.
    \item Découper en chunks.
    \item Ajouter les métadonnées : source, titre, section, date, langue, thème.
    \item Créer les embeddings.
    \item Sauvegarder l'index vectoriel.
\end{enumerate}

\subsection{Question}

\begin{enumerate}[leftmargin=1.2cm]
    \item Recevoir la question utilisateur.
    \item Créer l'embedding de la question.
    \item Rechercher les passages proches.
    \item Appliquer un reranking lexical.
    \item Construire une réponse structurée.
    \item Afficher les sources et la prudence juridique.
\end{enumerate}

\section{Modules réalisés}

\subsection{Assistant juridique}

Le module principal récupère les passages les plus pertinents, puis construit une réponse
avec les sections suivantes :

\begin{itemize}[leftmargin=1.2cm]
    \item réponse courte ;
    \item détails ;
    \item conséquences possibles ;
    \item procédure ;
    \item sources ;
    \item prudence juridique ;
    \item niveau de confiance.
\end{itemize}

\subsection{Calculateur d'amende et de points}

Le calculateur utilise le fichier :

\begin{lstlisting}
data/structured/infractions_maroc.csv
\end{lstlisting}

Il retourne le montant indicatif selon le délai de paiement, le nombre de points retirés,
la classe de l'infraction et la source officielle.

\subsection{Procédures}

Les procédures sont décrites dans :

\begin{lstlisting}
data/structured/procedures.json
\end{lstlisting}

Elles couvrent notamment :

\begin{itemize}[leftmargin=1.2cm]
    \item paiement d'une contravention ;
    \item déclaration d'un autre conducteur ;
    \item réclamation ;
    \item compréhension du retrait de points.
\end{itemize}

\subsection{Sources et vérification}

Chaque source est décrite dans :

\begin{lstlisting}
data/raw/sources_manifest.json
\end{lstlisting}

Ce manifeste permet de garder une traçabilité claire entre les réponses et les documents
institutionnels.

\section{Commandes principales}

\subsection{Construire l'index}

\begin{lstlisting}[language=bash]
python scripts/build_index.py
\end{lstlisting}

\subsection{Poser une question}

\begin{lstlisting}[language=bash]
python scripts/ask.py "Combien de points sont retirés pour un feu rouge ?"
\end{lstlisting}

\subsection{Calculer une amende}

\begin{lstlisting}[language=bash]
python scripts/calculate_fine.py --infraction "feu rouge" --delay 24h
\end{lstlisting}

\subsection{Lancer Streamlit}

\begin{lstlisting}[language=bash]
streamlit run app/streamlit_app.py
\end{lstlisting}

\section{Évaluation}

Le projet contient un petit jeu de questions d'évaluation :

\begin{lstlisting}
data/structured/evaluation_questions.csv
\end{lstlisting}

La commande :

\begin{lstlisting}[language=bash]
python scripts/evaluate.py
\end{lstlisting}

vérifie si la première source retrouvée correspond à la source attendue.

\section{Tests}

La suite de tests couvre :

\begin{itemize}[leftmargin=1.2cm]
    \item le découpage en chunks ;
    \item le calculateur d'amende ;
    \item les procédures ;
    \item l'indexation et la récupération de sources.
\end{itemize}

\begin{lstlisting}[language=bash]
pytest
\end{lstlisting}

\section{Limites}

\begin{itemize}[leftmargin=1.2cm]
    \item Le corpus versionné est une base initiale, pas une reproduction complète des textes officiels.
    \item Les PDF officiels doivent être téléchargés localement avec prudence.
    \item Le backend local \texttt{hashing} sert aux démonstrations ; un backend OpenAI ou FAISS est préférable en production.
    \item Les montants et points doivent être revérifiés lors de la mise à jour des sources.
\end{itemize}

\section{Améliorations possibles}

\begin{itemize}[leftmargin=1.2cm]
    \item Intégration complète des PDF SGG et NARSA.
    \item Extraction automatique des tableaux PDF.
    \item Ajout d'un reranker plus avancé.
    \item Passage vers FAISS ou ChromaDB.
    \item Ajout d'une interface bilingue français-arabe.
    \item Historique des conversations et export PDF de la réponse.
\end{itemize}

\section{Conclusion}

Tariqi Legal AI fournit une base complète et extensible pour un assistant RAG juridique
du code de la route marocain. Le projet combine recherche documentaire, citations,
données structurées, interface utilisateur et tests automatisés, tout en conservant une
prudence juridique claire.

\end{document}
